\chapter{Riesz の表現定理}
\label{ch:riesz-representation}

Riesz の表現定理は，連続関数に対する正の線形汎関数を，測度による積分として表す定理である．
第\ref{ch:disintegration}章では，Radon--Nikodym 微分から作った汎関数にこの定理を適用し，
各 $y$ に対応する確率測度を得る．ここでは compact 距離空間の場合を扱う．
$C(K)$ と $\lVert\cdot\rVert_\infty$ の記法は注意~\ref{rem:meas-compact-continuous}による．

\section{定理の主張と準備}

\begin{theorem}{Riesz--Markov--Kakutani の表現定理（compact 距離空間の場合）}{meas-riesz-representation}
  $K$ を空でない compact 距離空間とする．
  実線形写像 $L\colon C(K)\to\R$ が
  \[
    f\geq0\ \Longrightarrow\ L(f)\geq0,\qquad L(1)=1
  \]
  をみたすならば，ある唯一の $\rho\in\Prob(K)$ により
  \[
    L(f)=\int_Kf\,\d\rho\qquad(f\in C(K))
  \]
  と表せる．条件 $f\geq0\Rightarrow L(f)\geq0$ をみたす実線形写像を
  \textbf{正の線形汎関数}という．
\end{theorem}

証明は Rudin \emph{Real and Complex Analysis}（第3版）Theorem 2.14 の方針に従う．
まず開集合の測度を $L$ の値の上限で定め，Carathéodory の定理で Borel 測度を作る．
次に連続関数の積分が $L$ に一致することを示す．以下，$K,L$ は上の定理のとおりとする．

$f\leq g$ ならば $L(g)-L(f)=L(g-f)\geq0$ なので，$L$ は単調である．
特に $-\lVert f\rVert_\infty\leq f\leq\lVert f\rVert_\infty$ と $L(1)=1$ から
\[
  \abs{L(f)}\leq\lVert f\rVert_\infty
\]
を得る．

\begin{remark}{記法：開集合の内側で用いる連続関数}{rz-subordinate}
  開集合 $U\subset K$ と $g\in C(K)$ について，
  \[
    0\leq g\leq1,\qquad g=0\ \text{（$K\setminus U$ 上）}
  \]
  が成り立つとき，$g\prec U$ と書く．
\end{remark}

\begin{lemma}{開被覆に合わせた連続関数の分割}{rz-partition}
  $F\subset K$ を閉集合，$U_1,\dots,U_m\subset K$ を $F\subset U_1\cup\cdots\cup U_m$ をみたす開集合とし，
  $g\in C(K)$ は $0\leq g\leq1$ かつ $K\setminus F$ 上で $g=0$ をみたすとする．
  このとき $g_i\prec U_i$（$i=1,\dots,m$）で $g=g_1+\cdots+g_m$ となるものが存在する．
\end{lemma}

\begin{proof}
  $F=\emptyset$ なら $g_i:=0$ とすればよい．$F\neq\emptyset$ とする．
  各 $i$ について，$U_i\neq K$ なら $\phi_i(x):=d(x,K\setminus U_i)$，
  $U_i=K$ なら $\phi_i:=1$ とおく．
  補題~\ref{lem:meas-distance-to-set}より $\phi_i$ は連続で，$\phi_i(x)>0\iff x\in U_i$ である．
  $\Phi:=\sum_i\phi_i$ は compact 集合 $F$ 上で正なので
  $c:=\min_F\Phi>0$ である（注意~\ref{rem:meas-compact-facts}）．
  \[
    g_i:=\frac{g\,\phi_i}{\max\{\Phi,c\}}
  \]
  とおくと，$g_i$ は連続で非負であり，$U_i$ の外で $0$ である．
  $F$ 上では分母が $\Phi$ に等しく，$F$ の外では $g=0$ なので，$\sum_ig_i=g$ である．
  よって $0\leq g_i\leq g\leq1$ となり，$g_i\prec U_i$ である．
\end{proof}

この補題を $g\prec U$ に適用するには，$g$ を少し切り落とす．
$0<\varepsilon<1$ に対して
\[
  g_\varepsilon:=\max\{g-\varepsilon,0\},\qquad F_\varepsilon:=\{g\geq\varepsilon\}
\]
とおけば，$F_\varepsilon$ は $U$ に含まれる compact 集合であり，
$g_\varepsilon$ はその外で $0$ である．また $0\leq g-g_\varepsilon\leq\varepsilon$ なので
$L(g)-\varepsilon\leq L(g_\varepsilon)$ である．

\section{外測度の構成}

\begin{definition}{$L$ から定まる外測度}{rz-outer-measure}
  開集合 $U\subset K$ に対し
  \[
    \rho^*(U):=\sup\{L(g)\mid g\in C(K),\ g\prec U\}
  \]
  と定め，一般の $E\subset K$ に対し
  \[
    \rho^*(E):=\inf\{\rho^*(U)\mid U\text{ は開集合で }E\subset U\}
  \]
  と定める．
\end{definition}

開集合上で $\rho^*$ は単調なので，二つの定義は開集合上で一致する．
一般の集合についても単調性が成り立ち，$0\leq\rho^*(E)\leq1$，
$\rho^*(\emptyset)=0$，$\rho^*(K)=1$ である．

\begin{lemma}{$\rho^*$ は外測度である}{rz-is-outer-measure}
  $\rho^*$ は $K$ 上の外測度（定義~\ref{def:meas-outer-measure}）である．
\end{lemma}

\begin{proof}
  残る可算劣加法性を示す．まず $U=\bigcup_nU_n$ を開集合の合併とし，
  $g\prec U$，$0<\varepsilon<1$ を取る．前節の $F_\varepsilon$ は compact なので，
  相異なる有限個 $n_1,\dots,n_m$ で $F_\varepsilon\subset U_{n_1}\cup\cdots\cup U_{n_m}$ となる．
  補題~\ref{lem:rz-partition}で $g_\varepsilon=\sum_ig_i$，$g_i\prec U_{n_i}$ と分ければ
  \[
    L(g)-\varepsilon\leq L(g_\varepsilon)=\sum_{i=1}^mL(g_i)
    \leq\sum_{i=1}^m\rho^*(U_{n_i})\leq\sum_n\rho^*(U_n).
  \]
  $F_\varepsilon=\emptyset$ なら $g_\varepsilon=0$ なので，同じ評価が成り立つ．
  $\varepsilon\to0$ として $g$ について上限を取ると，$\rho^*(U)\leq\sum_n\rho^*(U_n)$ を得る．

  一般の $E=\bigcup_nE_n$ に対しては，$\varepsilon>0$ を取り，
  各 $n$ について開集合 $U_n\supset E_n$ を
  $\rho^*(U_n)\leq\rho^*(E_n)+\varepsilon2^{-n}$ となるように選ぶ．開集合での結果より
  \[
    \rho^*(E)\leq\rho^*\left(\bigcup_nU_n\right)\leq\sum_n\rho^*(U_n)
    \leq\sum_n\rho^*(E_n)+\varepsilon.
  \]
  $\varepsilon\to0$ として可算劣加法性を得る．
\end{proof}

\begin{lemma}{開集合の可測性}{rz-open-measurable}
  任意の開集合 $U\subset K$ は $\rho^*$-可測である．
\end{lemma}

\begin{proof}
  任意の $E\subset K$ について $\rho^*(E)\geq\rho^*(E\cap U)+\rho^*(E\setminus U)$ を示せばよい．

  \emph{$E=W$ が開集合の場合．}
  $0<\varepsilon<1$ とし，$g\prec W\cap U$ を
  $L(g)\geq\rho^*(W\cap U)-\varepsilon$ となるように取る．
  前節の $g_\varepsilon$，$F_\varepsilon$ を考えると，
  $W\setminus F_\varepsilon$ は $W\setminus U$ を含む開集合である．
  $k\prec W\setminus F_\varepsilon$ を $L(k)\geq\rho^*(W\setminus F_\varepsilon)-\varepsilon$ となるように取る．
  $g_\varepsilon$ と $k$ が正になる点は重ならず，$g_\varepsilon+k\prec W$ となるので，
  \[
    \begin{aligned}
      \rho^*(W)&\geq L(g_\varepsilon)+L(k)\\
      &\geq\rho^*(W\cap U)+\rho^*(W\setminus F_\varepsilon)-3\varepsilon\\
      &\geq\rho^*(W\cap U)+\rho^*(W\setminus U)-3\varepsilon.
    \end{aligned}
  \]
  $\varepsilon\to0$ として主張を得る．

  \emph{一般の場合．}
  $E\subset W$ となる開集合 $W$ について，前段と単調性より
  \[
    \rho^*(W)\geq\rho^*(W\cap U)+\rho^*(W\setminus U)
    \geq\rho^*(E\cap U)+\rho^*(E\setminus U).
  \]
  $W$ について下限を取ればよい．
\end{proof}

\section{表現定理の証明}

\begin{proof}[定理~\ref{thm:meas-riesz-representation}の証明]
  \textbf{1. 測度の構成．}
  補題~\ref{lem:rz-is-outer-measure}と Carathéodory の定理（定理~\ref{thm:meas-caratheodory}）より，
  $\rho^*$-可測集合の全体 $\mathcal{M}$ は $\sigma$-代数であり，$\rho^*$ は $\mathcal{M}$ 上の測度である．
  補題~\ref{lem:rz-open-measurable}より $\Borel(K)\subset\mathcal{M}$ なので，
  $\rho:=\rho^*|_{\Borel(K)}$ は $\rho(K)=1$ をみたす Borel 確率測度である．
  定義から，開集合 $U$ と $g\prec U$ について $L(g)\leq\rho(U)$ である．

  \textbf{2. 積分表示．}
  まず $f\in C(K)$，$0\leq f\leq1$ とする．$n\in\N$，$k=1,\dots,n$ に対して
  \[
    f_k:=\min\left\{\max\left\{f-\frac{k-1}n,0\right\},\frac1n\right\},\qquad
    U_k:=\left\{f>\frac{k-1}n\right\}
  \]
  とおく．これは $f$ を高さ $1/n$ の層に分ける操作であり，
  $f=\sum_{k=1}^nf_k$，$nf_k\prec U_k$ である．また，各点で
  \[
    f\leq\frac1n\sum_{k=1}^n\mathbf{1}_{U_k}\leq f+\frac1n
  \]
  が成り立つので，
  \[
    L(f)=\sum_{k=1}^nL(f_k)
    \leq\frac1n\sum_{k=1}^n\rho(U_k)
    \leq\int_Kf\,\d\rho+\frac1n.
  \]
  $n\to\infty$ として $L(f)\leq\int_Kf\,\d\rho$ を得る．
  $1-f$ にも適用すると，$L(1)=\rho(K)=1$ より逆向きの不等式も従う．
  一般の $f\in C(K)$ については，$m:=\lVert f\rVert_\infty>0$ ならば
  $g:=(f+m)/(2m)$ にこの等式を適用し，線形性から結論を得る．$m=0$ の場合は明らかである．

  \textbf{3. 一意性．}
  $\rho'\in\Prob(K)$ も同じ表示をみたすとする．
  空でない閉集合 $F$ について補題~\ref{lem:meas-closed-approx}の $\phi_k$ を取ると，
  有界収束定理（定理~\ref{thm:meas-bounded-convergence}）より
  \[
    \rho(F)=\lim_k\int_K\phi_k\,\d\rho=\lim_kL(\phi_k)
    =\lim_k\int_K\phi_k\,\d\rho'=\rho'(F).
  \]
  閉集合全体は $\Borel(K)$ を生成する $\pi$-系なので，
  一致定理（定理~\ref{thm:meas-uniqueness}）より $\rho=\rho'$ である．
\end{proof}
